\documentclass{situationist-preprint}

\title{Explanation and Prediction in Political Science}
\msid{POL-22-a}
\msfamily{POL-22}
\msdiscipline{Political Philosophy}
\mstype{Principal manuscript}
\mskeywords{political science; explanation; prediction; symmetry thesis; reflexivity; preference falsification; forecasting}

\begin{document}
\maketitle

\begin{abstract}
Can politics be a science? We argue that it can, in part: political science explains well but predicts poorly. Explanation and prediction come apart once explanation is understood as locating an outcome in a pattern of dependence rather than deriving it from laws and initial conditions. Political science has discovered many such patterns, such as the dependence of the number of parties on the electoral system. Prediction requires more: that the values of the relevant factors be known in advance, that no unmeasured factor intervene, and that the pattern itself remain stable once the prediction is made. Three features of political phenomena defeat these requirements: reflexivity, since agents respond to predictions and to the theories behind them; concealed preferences, which make large shifts invisible until they occur; and the multiplicity of interacting factors in any particular case.
\end{abstract}

\section{The question}

Behavioral political science aimed at general laws of political behavior tested by prediction; \citet{Popper1957} argued that the long-term course of history cannot be predicted, since it depends on the growth of knowledge, which cannot be predicted. The record of the discipline gives each side material. Political scientists failed, almost without exception, to predict the collapse of the communist regimes of Eastern Europe in 1989 \citep{Kuran1991}; and in large forecasting studies, experts' long-range political predictions did little better than simple extrapolation \citep{Tetlock2005}. Yet the discipline has well-confirmed explanatory generalizations.

We argue that the two observations are consistent: political science explains well but predicts poorly. The argument depends on the account of explanation defended in SCI-03, and it places political science alongside economics as treated in SCI-14, as a science whose models apply to actual cases only approximately.

\section{Explanation without symmetry}

On the deductive-nomological model, an explanation is an argument from laws and initial conditions to the explanandum, and every adequate explanation could have served, before the event, as a prediction \citep{HempelOppenheim1948}. If this \emph{symmetry thesis} were true, a science that predicted poorly would explain poorly too.

\begin{assumption}[Explanation as dependence]\label{a:dependence}
An explanation is an answer to a why-question that locates the explanandum in a pattern of dependence: it cites factors and a generalization relating them to the explanandum that would continue to hold under some changes in the factors.
\end{assumption}

Assumption~\ref{a:dependence} is the conclusion of SCI-03. On it, explanation does not require a law that guarantees the outcome, and the symmetry thesis fails: \citet{Scriven1959} observed that syphilis explains paresis although only a minority of those with untreated syphilis develop it, so that the explanation would not have licensed the prediction.

\begin{definition}\label{d:prediction}
A generalization \emph{supports prediction} of an outcome in a case if, before the outcome, (i) the values of the factors it cites are known, (ii) no factor it does not cite will change the outcome, and (iii) the generalization will continue to hold in the case once the prediction is made.
\end{definition}

\section{The asymmetry}

\begin{proposition}[Political science explains well]\label{p:explains}
Political science possesses stable patterns of dependence that explain particular political outcomes.
\end{proposition}

\begin{proof}
By Assumption~\ref{a:dependence}, it suffices to exhibit generalizations that relate political outcomes to factors and hold under changes in those factors. Duverger's law relates the number of competitive parties to the electoral system: plurality rule in single-member districts tends toward two-party competition, through the mechanical under-representation of third parties and voters' reluctance to waste their votes \citep{Duverger1954}. The theory of collective action relates the provision of public goods by groups to their size and to the availability of selective incentives \citep{Olson1965}. Both generalizations have exceptions, but both hold across many changes in the cited factors and identify the mechanisms through which they operate, which is what an explanation by dependence requires. As with the idealized models of economics (SCI-14), they apply to a particular case to the extent that the factors they omit do not make a difference there.
\end{proof}

\begin{proposition}[Political science predicts poorly]\label{p:predicts}
For the outcomes that matter most in politics, such as changes of regime, shifts of party systems and the outbreak of conflict, the conditions of Definition~\ref{d:prediction} generally fail.
\end{proposition}

\begin{proof}
Condition (iii) fails because political agents respond to predictions and to the theories that generate them. A forecast of a party's collapse affects donors and voters; a theory of electoral behavior is used by campaigns to change behavior; a model of policy effects changes once agents anticipate the policy, the point of the Lucas critique in macroeconomics \citep{Lucas1976}, which applies whenever the agents studied can learn the model. Condition (i) fails because the factors on which large shifts depend include concealed preferences: under repression citizens falsify their preferences in public, so that the distribution of private opposition, and the thresholds at which individuals would join a protest, are unknown until a cascade reveals them \citep{Kuran1991}. Condition (ii) fails because in any particular case many factors interact, any of which can override the effect of those cited, so that a generalization true in most cases offers no assurance in this one. Each failure is compatible with the explanatory generalizations of Proposition~\ref{p:explains}: after the event, the cascade, the anticipation or the intervening factor can be identified and the outcome located in the pattern.
\end{proof}

\begin{theorem}\label{t:main}
Politics can be a science in part: political science explains well but predicts poorly.
\end{theorem}

\begin{proof}
Proposition~\ref{p:explains} gives the explanatory success; Proposition~\ref{p:predicts} gives the predictive failure; and by Assumption~\ref{a:dependence} the two are consistent, since explanation does not require the conditions that prediction requires. Forecasting studies agree: forecasters do best on short horizons and narrow questions, where the conditions of Definition~\ref{d:prediction} nearly hold, as in election forecasts from polls, and worst on long horizons \citep{Tetlock2005}.
\end{proof}

\section{Objections}

\begin{objection}[Explanation without prediction is storytelling]
A discipline that can explain any outcome after the fact but predict none has no way of testing its explanations. Its explanations are stories.
\end{objection}

\begin{reply}
Explanatory generalizations are tested across many cases, not by forecasts of particular events. Duverger's law is tested by comparing party systems across electoral rules and by the effects of electoral reforms; claims about tendencies in populations can be confirmed when individual outcomes cannot be predicted.
\end{reply}

\begin{objection}[Poor prediction shows immaturity]
The predictive weakness of political science reflects its youth and its data, and will be overcome.
\end{objection}

\begin{reply}
Better data will relax condition (ii) in some cases. Conditions (i) and (iii) do not depend on the quantity of data: concealed preferences are concealed from data collection, and reflexivity grows with the discipline's influence. The limit is structural, not a matter of maturity.
\end{reply}

\section{Conclusion}

Political science has found patterns of dependence that explain political outcomes, and these satisfy the conditions of explanation without satisfying those of prediction. Agents respond to predictions, conceal the preferences on which large shifts depend, and act in cases where many factors interact. Politics can be a science of explanation, and only in a limited way a science of prediction.

\begin{thebibliography}{99}\small
\bibitem[Duverger(1954)]{Duverger1954} Duverger, M. (1954). \emph{Political Parties: Their Organization and Activity in the Modern State} (B. North and R. North, trans.). Methuen.
\bibitem[Hempel and Oppenheim(1948)]{HempelOppenheim1948} Hempel, C. G., and Oppenheim, P. (1948). Studies in the logic of explanation. \emph{Philosophy of Science}, 15(2), 135--175.
\bibitem[Kuran(1991)]{Kuran1991} Kuran, T. (1991). Now out of never: The element of surprise in the East European revolution of 1989. \emph{World Politics}, 44(1), 7--48.
\bibitem[Lucas(1976)]{Lucas1976} Lucas, R. E. (1976). Econometric policy evaluation: A critique. \emph{Carnegie-Rochester Conference Series on Public Policy}, 1, 19--46.
\bibitem[Olson(1965)]{Olson1965} Olson, M. (1965). \emph{The Logic of Collective Action: Public Goods and the Theory of Groups}. Harvard University Press.
\bibitem[Popper(1957)]{Popper1957} Popper, K. R. (1957). \emph{The Poverty of Historicism}. Routledge \& Kegan Paul.
\bibitem[Scriven(1959)]{Scriven1959} Scriven, M. (1959). Explanation and prediction in evolutionary theory. \emph{Science}, 130(3374), 477--482.
\bibitem[Tetlock(2005)]{Tetlock2005} Tetlock, P. E. (2005). \emph{Expert Political Judgment: How Good Is It? How Can We Know?} Princeton University Press.
\end{thebibliography}

\end{document}
